\hypertarget{chapter-01-page-013}{}
\subsection{Fourier 系数与级数}
\label{sec:ch1-fourier-coefficients-series}

在(一个区间上的)$\mathbb R$ 上定义的函数 $f$ 用形如
\begin{equation}
\label{eq:ch1-trigonometric-series}
 f(x)=\sum_{k=0}^{\infty}a_k\cos(kx)+b_k\sin(kx)
\end{equation}
的三角级数表示这一问题, 在用分离变量法求解偏微分方程时自然出现. J. Fourier 正是由此接触到这个问题, 并将其著作 \emph{Th\'eorie Analytique de la Chaleur} 的大部分篇幅用于研究它(该书出版于 1822 年, 其结果于 1807 年首次提交给 Institute de France). 更早以前, 在 18 世纪中叶, Daniel Bernoulli 试图求解弦振动问题时已经提出了这一问题, 而系数公式则出现于 L. Euler 1777 年的一篇论文中.

\eqref{eq:ch1-trigonometric-series} 右端是周期为 $2\pi$ 的周期函数, 因而 $f$ 也必须具有这一性质. 所以只需在一个长度为 $2\pi$ 的区间上考虑 $f$. 利用 Euler 恒等式 $e^{ikx}=\cos(kx)+i\sin(kx)$, 可将 \eqref{eq:ch1-trigonometric-series} 中的函数 $\sin(kx)$ 和 $\cos(kx)$ 换成 $\{e^{ikx}:k\in\mathbb Z\}$; 从现在起我们都这样做. 此外, 我们将考虑周期为 $1$ 而不是 $2\pi$ 的函数, 因而把函数系改为 $\{e^{2\pi ikx}:k\in\mathbb Z\}$. 这样, 我们的问题就转化为研究 $f$ 的表示
\begin{equation}
\label{eq:ch1-complex-fourier-series}
 f(x)=\sum_{k=-\infty}^{\infty}c_ke^{2\pi ikx}.
\end{equation}

\hypertarget{chapter-01-page-014}{}
例如, 若假设该级数一致收敛, 则乘以 $e^{-2\pi imx}$ 并在 $(0,1)$ 上逐项积分, 得到
\[
 c_m=\int_0^1f(x)e^{-2\pi imx}\,dx,
\]
这是因为正交关系
\begin{equation}
\label{eq:ch1-exponential-orthogonality}
 \int_0^1e^{2\pi ikx}e^{-2\pi imx}\,dx=
 \begin{cases}
 0,&k\ne m,\\
 1,&k=m.
 \end{cases}
\end{equation}

以 $\mathbb T$ 表示实数模 $1$ 的加法群(即 $\mathbb R/\mathbb Z$), 也就是一维环面. 它也可与单位圆 $S^1$ 等同. 说一个函数定义在 $\mathbb T$ 上, 等价于说它定义在 $\mathbb R$ 上且周期为 $1$. 对每个函数 $f\in L^1(\mathbb T)$, 赋予其 Fourier 系数序列 $\{\widehat f(k)\}$, 定义为
\begin{equation}
\label{eq:ch1-fourier-coefficient}
 \widehat f(k)=\int_0^1f(x)e^{-2\pi ikx}\,dx.
\end{equation}
以这些数为系数的三角级数
\begin{equation}
\label{eq:ch1-fourier-series}
 \sum_{k=-\infty}^{\infty}\widehat f(k)e^{2\pi ikx}
\end{equation}
称为 $f$ 的 Fourier 级数.

现在的问题是确定级数 \eqref{eq:ch1-fourier-series} 在何时以及按何种意义表示函数 $f$.

